\documentclass[12pt,a4paper,oneside]{article}
\usepackage[brazil]{babel}
\usepackage[utf8]{inputenc}
\usepackage{olymp}

\setcounter{problem}{3}

\begin{document}

\begin{problem}{Campo de minhocas}{minhocas}{2}

Minhocas são muito importantes para a agricultura e como insumo para produção de ração animal. A Organização para Bioengenharia de Minhocas (OBM) é uma entidade não governamental que promove o aumento da produção, utilização e exportação de minhocas.\\[-18pt]

Uma das atividades promovidas pela OBM é a manutenção de uma fazenda experimental para pesquisa de novas tecnologias de criação de minhocas. Na fazenda, a área destinada às pesquisas é de formato retangular, dividida em células quadradas de mesmo tamanho. %Em cada célula é criada apenas uma espécie de minhoca.
%As células são utilizadas para testar os efeitos, sobre a produção de minhocas, de variações de espécies de minhoca, de tipos de terra, de adubo, de umidade, etc.
Os pesquisadores da OBM mantêm um acompanhamento constante do desenvolvimento das minhocas em cada célula, e têm uma estimativa extremamente precisa da produtividade de cada uma das células.\\[-18pt]

\begin{center}
\begin{tabular}{|c|c|c|c|}
\hline
81 & 28 & 240 & 10 \\
\hline
40 & 10 & 100 & 240 \\
\hline
20 & 180 & 110 & 35 \\
\hline
\end{tabular}
\end{center}
\vspace{-6pt}

Um pesquisador da OBM inventou e construiu uma máquina colhedeira de minhocas, e quer testá-la na fazenda. A máquina tem a largura de uma célula, e em uma passada pelo terreno de uma célula colhe todas as minhocas dessa célula, separando-as, limpando-as e empacotando-as.
%Ou seja, a máquina eliminara uma das etapas mais intensivas de mão de obra no processo de produção de minhocas.
A máquina, porém, ainda está em desenvolvimento e %tem uma restrição: não faz curvas, podendo movimentar-se somente em linha reta.
não faz curvas, move-se apenas em linha reta.\\[-18pt]

Decidiu-se %então 
que seria efetuado um teste com a máquina, de forma a colher o maior número possível de minhocas em uma única passada, em linha reta, de lado a lado do campo de minhocas. Ou seja, a máquina deve colher todas as minhocas de uma `coluna' ou de uma `linha' de células do campo de minhocas (a linha ou coluna cuja soma das produtividades %esperadas 
das células é a maior possível).\\[-18pt]

Escreva um programa que, fornecido o mapa do campo de minhocas
%, descrevendo a produtividade
com a produtividade
estimada em cada célula,
calcule o %número esperado total
total esperado
de minhocas a serem colhidas pela máquina %durante o teste, conforme descrito acima.
no teste.

%\Notes
%
%Observações, se tiver.

\InputFile

A primeira linha da entrada contém dois %números
inteiros $N$ e $M$, representando respectivamente o número de linhas e %o número 
de colunas de células existentes no campo experimental de minhocas. Cada uma das $N$ linhas seguintes contém $M$ inteiros, representando as produtividades estimadas das células correspondentes a uma linha do campo de minhocas.
Restrições: $1 \leq N \leq 100; 1 \leq M \leq 100$.

\OutputFile

%A saída deve ser composta por uma única linha contendo um inteiro, indicando o número esperado total de minhocas a serem colhidas pela máquina durante o teste.
A saída deve ser um inteiro informando o total esperado de minhocas colhidas pela máquina.

% \Notes
% É garantido que sempre existe uma divisão que satisfaz as condições dos reinos.

\Example

\begin{example}
\exmp{
3 4
81 28 240 10
40 10 100 240
20 180 110 35
}{
450
}%
\end{example}

\begin{example}
\exmp{
3 1
110
0
100
}{
210
}%
\end{example}

\small \textit{Obs.: questão da OBI 2005} 

%Comentários sobre o exemplo, se necessário.

\end{problem}

\end{document}
